\documentclass[10pt,twocolumn]{article}
\usepackage[T1]{fontenc}
\usepackage[utf8]{inputenc}
\usepackage{lmodern}
\usepackage[margin=1.8cm,top=2.4cm,bottom=2cm,columnsep=0.6cm]{geometry}
\usepackage{fancyhdr}
\usepackage{titlesec}
\usepackage{booktabs}
\usepackage{amsmath,amssymb,amsfonts}
\usepackage{microtype}
\usepackage{xcolor}
\usepackage{mdframed}
\usepackage{parskip}
\usepackage{enumitem}
\usepackage{hyperref}

\definecolor{rulered}{RGB}{180,30,30}
\definecolor{boxgray}{RGB}{245,245,245}
\definecolor{keywordgray}{RGB}{100,100,100}

\pagestyle{fancy}
\fancyhf{}
\fancyhead[L]{\textsc{\small Claude Mag}}
\fancyhead[C]{\small\textit{Machine Learning Research Digest}}
\fancyhead[R]{\small 2026-09-21}
\fancyfoot[C]{\small\thepage}
\renewcommand{\headrulewidth}{0.8pt}

\titleformat{\section}{\normalsize\bfseries\color{rulered}}{}{0pt}{}%
  [\vspace{-4pt}\color{rulered}\rule{\columnwidth}{0.4pt}]
\titlespacing{\section}{0pt}{8pt}{4pt}

\hypersetup{colorlinks=true,urlcolor=rulered,linkcolor=black}

\begin{document}
\setlength{\parindent}{0pt}
\setlength{\parskip}{4pt}

{\small\color{keywordgray}\textbf{Keywords:}
  hierarchical robot planning \textbullet{}
  vision-language models \textbullet{}
  agentic distillation \textbullet{}
  stage-transition detection}\par\vspace{4pt}

{\LARGE\bfseries\raggedright
  StageGuard}\par\vspace{4pt}

{\large\itshape\raggedright
  Three-Layer Agentic Distillation Teaches a Compact VLM to Decide
  When Robots Should Switch Subtasks---at Real-Time Speed}\par\vspace{6pt}

{\small\textbf{By} Wenbo Cui, Jing Xu, Rui Zhao, Yunchao Liu, Chaobin Ye,
  Xinyu Zhang \& Mingkui Tan (South China University of Technology;
  Tencent Robotics~X)}\par\vspace{2pt}

{\small\color{keywordgray}arXiv:2609.20791 \textbullet{} Submitted September 17, 2026}
\par\vspace{2pt}\noindent{\color{rulered}\rule{\columnwidth}{0.6pt}}\vspace{6pt}

The bottleneck in long-horizon robot planning is not the skills themselves
but knowing when to stop one and start the next. \textit{StageGuard} distils
structured three-layer reasoning from a large cloud VLM into a compact local
model that runs at 11~Hz, achieves F1\,=\,88.4 on stage-transition
detection---7.2 points above the best prior method---and improves full-task
success rate by 9.7 percentage points on BEHAVIOR-1K.

\begin{mdframed}[backgroundcolor=boxgray,linecolor=rulered,linewidth=1pt,
    innerleftmargin=6pt,innerrightmargin=6pt,
    innertopmargin=6pt,innerbottommargin=6pt]
\textbf{\small AT A GLANCE}\par\vspace{4pt}
\begin{description}[leftmargin=0pt,labelsep=4pt,itemsep=2pt,parsep=0pt,
    font=\small\bfseries]
  \item[Problem:] Stage-transition decisions in hierarchical robot
    planning require human-designed completion checkers that fail under
    perceptual noise; large VLMs can reason about completion but are
    misaligned with subtask criteria and too slow for real-time control.
  \item[Why it matters:] Errors in stage transitions cascade: a premature
    transition moves to the next skill before the current one is complete,
    and a missed transition stalls the robot indefinitely.
  \item[Method:] A teacher VLM generates structured three-layer traces
    (state, task, decision) from demonstration trajectories; a compact
    student VLM is fine-tuned on these traces and deployed on-robot.
  \item[Result:] F1\,=\,88.4 (sim), 87~ms latency, 67.1\% task success
    rate on BEHAVIOR-1K; validated on a real Franka Panda.
  \item[Evidence:] Comparison with rule-based checkers, zero-shot GPT-4V,
    and fine-tuned VLM without structured reasoning; ablations over
    each reasoning layer and observation history length.
\end{description}
\end{mdframed}

\section{The Big Picture}

Hierarchical robot planners decompose complex tasks (e.g.,
``load the dishwasher and close it'') into a sequence of subtasks
(grasp plate, place plate, grasp glass, \ldots, close door), each
executed by a dedicated low-level skill policy. The stage-transition
monitor decides when to terminate the current skill and invoke the next.

This is a harder problem than it appears. Low-level skill policies
produce smooth trajectories that may overshoot, stall, or partially
succeed. The robot may not be able to observe all relevant state directly.
Transition criteria that are trivially stated in natural language
(``the drawer is fully closed'') require fine-grained spatial reasoning
under occlusion and varied lighting.

Current solutions are hand-crafted rule sets: threshold on end-effector
position, contact detection, object bounding-box IoU. These rules are
brittle, task-specific, and require re-engineering for each new subtask.

\section{Why Existing Approaches Fall Short}

\textbf{Rule-based completion checkers} must be designed per subtask by
a robotics engineer, do not generalise to novel objects or environments,
and fail under sensor noise. Their F1 on BEHAVIOR-1K is only 71.3.

\textbf{Cloud VLMs (GPT-4V, Qwen-VL-Max)} can interpret diverse scene
images but have three problems for this use case: (1) their query--answer
format is not calibrated to subtask-completion criteria specifically;
(2) they operate at 1--3 seconds per query vs.\ 10\,Hz control loops;
(3) they require cloud access, making them unsuitable for offline or
safety-critical deployments.

\textbf{Fine-tuned VLMs without structured reasoning} achieve F1\,=\,81.6
but have a flat, one-layer reasoning process: they output binary decisions
without decomposing state observations, task semantics, or decision
rationales separately, limiting their robustness.

\section{Core Mechanism}

The teacher VLM processes each timestep observation through three
sequential prompting stages that build a structured trace
$(z_{\text{state}},\,z_{\text{task}},\,z_{\text{decision}},\,d_t)$:

\textbf{State layer.} Given observation history $o_{t-H:t}$, the
teacher extracts a structured state summary: current object poses,
grasp contacts, spatial relations, and completion indicators (e.g.,
``cup rim 2~cm above shelf surface, no contact'').

\textbf{Task layer.} Given $z_{\text{state}}$ and the current/next
subtask pair $(T_i, T_{i+1})$, the teacher maps state facts to subtask
completion criteria: which aspects of $z_{\text{state}}$ determine
whether $T_i$ is done, and what preconditions $T_{i+1}$ requires.

\textbf{Decision layer.} Given $z_{\text{task}}$, the teacher issues
a binary transition verdict $d_t \in \{0,1\}$ with a natural-language
rationale $z_{\text{decision}}$.

A compact local student VLM (3B parameters) is fine-tuned on
$(o_{t-H:t}, T_i, T_{i+1}) \to (z_{\text{decision}}, d_t)$ pairs
generated by running the teacher on all demonstration trajectories.

\section{Technical Formulation}

Formally, the monitor observes $o_t = (I_t, \mathbf{q}_t, \mathbf{p}_t)$
(RGB image, joint configuration, end-effector pose) with history window
$H$. The transition decision:
\[
d_t = f(o_{t-H:t},\, T_i,\, T_{i+1}) \in \{0,1\}.
\]

The student produces a soft confidence from its last-layer hidden state:
\[
c_t = \sigma\!\left(w^\top h_t\right), \quad h_t = \text{Student}(o_{t-H:t}, T_i, T_{i+1}),
\]
triggering a transition when $c_t \geq \tau$ (set by validation).

The student training loss combines binary cross-entropy on the decision
and a token-level alignment loss on the compressed explanation:
\[
\mathcal{L}_S = \mathcal{L}_{\text{CE}}(d_t, \hat{d}_t)
              + \lambda\,\mathcal{L}_{\text{align}}(z_{\text{decision}},\hat{z}_{\text{decision}}).
\]
The alignment loss is a token-level cross-entropy weighted by teacher
confidence:
\[
\mathcal{L}_{\text{align}} = -\sum_j \hat{p}_j \log p_j^{\text{student}},
\]
where $\hat{p}_j$ are teacher token probabilities for explanation
token $j$.

\section{Learning or Inference Procedure}

\begin{enumerate}[itemsep=2pt,parsep=0pt]
  \item Collect human demonstration trajectories for target task set.
  \item Run cloud teacher VLM (offline) to produce structured
    traces for all trajectory timesteps.
  \item Fine-tune compact student VLM on trace-decision pairs.
  \item Deploy student on-robot: at each 10\,Hz control step, pass
    recent observation and subtask context; trigger transition if
    $c_t \geq \tau$.
  \item Log student explanation text for safety records and debugging.
\end{enumerate}

Training uses $\lambda = 0.3$, history $H=5$, threshold $\tau=0.55$.
Teacher traces are generated once and cached; student fine-tuning takes
approximately 4 hours on a single A100.

\section{What the Guarantee Says}

StageGuard provides no formal guarantees. Empirically, the student's F1
is within 2.3 points of the teacher's F1 on held-out BEHAVIOR-1K
trajectories, while inference latency drops from 1,820~ms (cloud VLM)
to 87~ms (local student), enabling real-time monitoring at 11~Hz.
Real-robot experiments on five Franka Panda tasks confirm the
improvements transfer beyond simulation.

\section{Experimental Findings}

\begin{center}
\small
\begin{tabular}{lccc}
\toprule
Method & F1 & Lat.\ (ms) & Task SR \\
\midrule
Rule-based & 71.3 & 2 & 48.2 \\
GPT-4V zero-shot & 79.1 & 1820 & 52.7 \\
VLM FT (no reasoning) & 81.6 & 91 & 57.4 \\
\textbf{StageGuard} & \textbf{88.4} & \textbf{87} & \textbf{67.1} \\
\bottomrule
\end{tabular}
\end{center}

Task SR = full long-horizon episode success rate on BEHAVIOR-1K simulation.
StageGuard improves task success by 9.7 points over the best prior method
(VLM FT) at essentially the same inference latency.

\section{Ablations and Interpretation}

\textbf{Removing state layer:} F1 drops 4.1 points; spatial
relationship reasoning degrades without the structured state summary.

\textbf{Removing task layer:} F1 drops 3.6 points; subtask completion
criteria are not grounded without explicit task-plan mapping.

\textbf{Unstructured distillation:} F1 = 83.9 (vs.\ 88.4); unstructured
teacher traces dilute relevant decision factors with spurious reasoning.

\textbf{History window $H$:} $H=5$ is optimal; $H=1$ loses 4.8 F1;
$H=10$ gains 0.4 F1 at double the latency.

\textbf{Student size:} 3B achieves full StageGuard F1; 1B loses 3.2 F1;
7B adds only 0.5 F1 at 3$\times$ latency.

\section*{Reference}

Wenbo Cui, Jing Xu, Rui Zhao, Yunchao Liu, Chaobin Ye, Xinyu Zhang,
Mingkui Tan.
\textbf{StageGuard: Learning Stage Transitions for Long-Horizon Robot Tasks
via Agentic Distillation.}
arXiv:2609.20791, September 2026.
\url{https://arxiv.org/abs/2609.20791}

\end{document}
